\documentclass[tikz, border=3mm]{standalone}
\usepackage{tikz}
\usetikzlibrary{calc, positioning, quotes, angles}

\begin{document}
\begin{tikzpicture}[>=stealth, scale=1.5]

% --- Definitions ---
% We define two different angles to match the drawing:
% 1. \theta: The angle shown for the solid lines.
% 2. \waveangle: The angle for the dashed lines (wavefronts).
\def\theta{35}      % Angle \theta as SHOWN (from vertical)
\def\waveangle{55} % Angle for DASHED lines (from vertical)
\def\d{3}          % Distance d

% --- Coordinates ---
\coordinate (r1) at (0, 0);
\coordinate (r0) at (0, \d);
\coordinate (P) at (7, 5); % Point P

% --- Z-Axis ---
\draw [->, thick] (0, -1) -- (0, \d+2) node[above] {$z$};
\fill (r0) circle (2pt) node [left=2pt] {$r_0$};
\fill (r1) circle (2pt) node [left=2pt] {$r_1$};

% --- Distance d ---
\draw [<->, red, thick] (r0) -- (r1) node [midway, left=2pt] {$d$};

% --- Solid Lines & Angle \theta ---
% These are the two solid lines.
\draw [thick] (r0) -- ++(90-\theta:9) coordinate (solid0);
\draw [thick] (r1) -- ++(90-\theta:9);
% Draw the angle \theta relative to the solid line, as in the image.
\coordinate (z_ref) at (0, \d+1.5); % Reference point on z-axis
\pic [draw, angle_radius=1cm, "$\theta$"] {angle = z_ref--r0--solid0};

% --- Dashed Lines (Wavefronts) ---
\def\tika{90-\waveangle} % Ti*k*Z angle for dashed lines
\def\tikp{-\waveangle}   % Ti*k*Z angle for perpendicular
% Draw the reference dashed line passing through r0
\draw [dashed] (r0) -- ++(\tika:8) coordinate (wave0); 
% Draw parallel dashed lines
\foreach \i in {1, 2, 3} {
    \draw [dashed] ($(r0)+(\tikp:\i*0.6)$) -- ++(\tika:9);
}
\foreach \i in {1} {
    \draw [dashed] ($(r0)+(\tikp:-\i*0.6)$) -- ++(\tika:9);
}

% --- Kappa Vector ---
% Positioned relative to the diagram, perpendicular to dashed lines.
\coordinate (kappa_start) at (3, 2);
\draw [->, red, thick] (kappa_start) -- ++(180+\tikp:1.5) node [midway, right=2pt] {$\kappa$};

% --- Path Difference ---
% Project r1 onto the reference dashed line (wave0).
\coordinate (p) at ($(r0)!(r1)!(wave0)$);
% Draw the blue line and label it "d cos(theta)" as in the image,
% even though its geometric length is d*cos(\waveangle).
\draw [->, blue, thick] (r1) -- (p) node [midway, below left, xshift=-2pt] {$d \cos \theta$};

% --- Point P ---
\fill (P) circle (2pt) node [right=2pt] {$P$};
\draw [dashed] (P) circle (0.5);
\draw [dashed] (P) circle (0.8);

\end{tikzpicture}
\end{document}